% sudodoc-elements.tex — Document Element Templates
% Tables, lists, code, blockquotes, figures, descriptions
% Compile: lualatex sudodoc-elements.tex (run TWICE for overlays)
\documentclass[12pt]{article}

% ─── Packages ───
\usepackage{fontspec}
\usepackage{geometry}
\usepackage{xcolor}
\usepackage{tikz}
\usepackage{graphicx}
\usepackage{enumitem}      % Load BEFORE titlesec to avoid conflict
\usepackage{titlesec}
\usepackage{titletoc}
\usepackage{parskip}
\usepackage{calc}
\usepackage{fancyhdr}
\usepackage{lastpage}
\usepackage{booktabs}
\usepackage{listings}
\usepackage{caption}
\usepackage{colortbl}
\usepackage{multirow}
\usepackage{tcolorbox}
\usepackage{amsmath}
\usepackage{amssymb}

% ─── Fonts ───
\setsansfont{Latin Modern Sans}
\setmainfont{Latin Modern Roman}
\setmonofont{Latin Modern Mono}

% ─── Colors (matching title page + TOC design) ───
\definecolor{titleDark}{HTML}{1A1A2E}
\definecolor{titleAccent}{HTML}{16213E}
\definecolor{titleRule}{HTML}{0F3460}
\definecolor{titleLight}{HTML}{E8E8E8}
\definecolor{titleMid}{HTML}{888888}

% Element-specific colors
\definecolor{tableHead}{HTML}{1A1A2E}
\definecolor{tableHeadText}{HTML}{FFFFFF}
\definecolor{tableAlt}{HTML}{F0F4F8}
\definecolor{tableBorder}{HTML}{D0D8E0}
\definecolor{codeBg}{HTML}{F5F7FA}
\definecolor{codeFrame}{HTML}{D0D8E0}
\definecolor{codeKeyword}{HTML}{0F3460}
\definecolor{codeString}{HTML}{2E7D32}
\definecolor{codeComment}{HTML}{90A4AE}
\definecolor{quoteBg}{HTML}{F8F4F0}
\definecolor{quoteBorder}{HTML}{0F3460}
\definecolor{quoteText}{HTML}{3E3E3E}
\definecolor{bulletColor}{HTML}{0F3460}
\definecolor{enumColor}{HTML}{0F3460}
\definecolor{descTerm}{HTML}{0F3460}
\definecolor{descText}{HTML}{3E3E3E}
\definecolor{captionLabel}{HTML}{0F3460}
\definecolor{captionText}{HTML}{666666}

% ─── Geometry ───
\geometry{
  paper=a4paper,
  top=30mm,
  bottom=30mm,
  left=30mm,
  right=25mm,
  headheight=15pt,
}

% ══════════════════════════════════════════════════════════
% SECTIONING (from sudodoc-toc)
% ══════════════════════════════════════════════════════════
\titleformat{\section}
  {\normalfont\Large\bfseries\color{titleDark}}
  {}{0pt}
  {\vspace{6pt}%
   \begin{tikzpicture}[baseline=(title.base)]
     \node[anchor=west, inner sep=0pt] (title) {\insertsection};
     \draw[titleRule, line width=1.2pt]
       ([yshift=-3pt]title.south west) -- ([yshift=-3pt]title.south east);
   \end{tikzpicture}%
  }
  [\vspace{2pt}]

\titleformat{\subsection}
  {\normalfont\large\bfseries\color{titleRule}}
  {}{0pt}
  {\hspace{4pt}%
   \tikz[baseline=(num.base)]{
     \node[fill=titleRule, text=white, rounded corners=2pt,
           inner xsep=6pt, inner ysep=2pt,
           font=\footnotesize\bfseries] (num)
       {\thesubsection};
   }%
   \hspace{6pt}\insertsubsection%
  }
  [\vspace{2pt}]

\titleformat{\subsubsection}
  {\normalfont\normalsize\bfseries\color{titleMid}}
  {}{0pt}
  {\hspace{14pt}{\color{titleRule}\textbullet}\hspace{8pt}\insertsubsubsection}
  [\vspace{1pt}]

\titlespacing{\section}{0pt}{24pt plus 4pt minus 2pt}{12pt plus 2pt minus 1pt}
\titlespacing{\subsection}{0pt}{18pt plus 3pt minus 1pt}{8pt plus 1pt minus 1pt}
\titlespacing{\subsubsection}{0pt}{12pt plus 2pt minus 1pt}{6pt}

% ══════════════════════════════════════════════════════════
% TABLES: booktabs + alternating rows
% ══════════════════════════════════════════════════════════

% ── Alternating row colors ──
\newcommand{\rowcolorhead}{\rowcolor{tableHead}}   % dark header
\newcommand{\rowcoloralt}{\rowcolor{tableAlt}}      % light alt row
\arrayrulecolor{tableBorder}

% ── Caption style (simple) ──
\captionsetup{
  tableposition=top,
  font=small,
  labelfont=bf,
  labelsep=period,
  skip=8pt,
  belowskip=6pt,
}

% ══════════════════════════════════════════════════════════
% LISTS: Custom bullets, numbers, descriptions
% ══════════════════════════════════════════════════════════

% ── Itemize: styled bullets (no TikZ to avoid enumitem conflict) ──
\setitemize{
  itemsep=6pt,
  parsep=2pt,
  topsep=8pt,
  partopsep=4pt,
  leftmargin=18pt,
  labelsep=8pt,
}
% Level 1: filled triangle
\renewcommand{\labelitemi}{\textcolor{bulletColor}{\small$\blacktriangleright$}}
% Level 2: filled circle
\renewcommand{\labelitemii}{\textcolor{titleRule}{\raisebox{0.2ex}{\tiny$\blacksquare$}}}
% Level 3: open square
\renewcommand{\labelitemiii}{\textcolor{titleMid}{\raisebox{0.1ex}{\scriptsize$\square$}}}

% ── Enumerate: styled number badges (no TikZ to avoid enumitem conflict) ──
\setenumerate{
  itemsep=6pt,
  parsep=2pt,
  topsep=8pt,
  partopsep=4pt,
  leftmargin=24pt,
  labelsep=8pt,
}
% Level 1: number in colored badge
\renewcommand{\labelenumi}{%
  \colorbox{enumColor}{\textcolor{white}{\footnotesize\bfseries\hspace{4pt}\theenumi\hspace{4pt}}}%
}
% Level 2: letter in gray badge
\renewcommand{\labelenumii}{%
  \colorbox{titleMid}{\textcolor{white}{\footnotesize\bfseries\hspace{4pt}\theenumii\hspace{4pt}}}%
}

% ── Description ──
\setdescription{
  itemsep=8pt,
  parsep=2pt,
  topsep=8pt,
  leftmargin=0pt,
  labelsep=12pt,
  font={\bfseries\color{descTerm}},
}

% ══════════════════════════════════════════════════════════
% CODE LISTINGS: Syntax highlighting
% ══════════════════════════════════════════════════════════
\lstdefinestyle{sudodoc}{
  backgroundcolor=\color{codeBg},
  basicstyle=\ttfamily\small\color{titleDark},
  keywordstyle=\bfseries\color{codeKeyword},
  stringstyle=\color{codeString},
  commentstyle=\itshape\color{codeComment},
  numberstyle=\tiny\color{titleMid},
  numbers=left,
  numbersep=10pt,
  frame=l,
  framerule=2pt,
  rulecolor=\color{codeKeyword},
  xleftmargin=16pt,
  framexleftmargin=16pt,
  breaklines=true,
  breakatwhitespace=true,
  tabsize=4,
  showstringspaces=false,
  captionpos=b,
  aboveskip=12pt,
  belowskip=6pt,
  columns=fullflexible,
  keepspaces=true,
}
\lstset{style=sudodoc}

% ── Code caption style ──
\captionsetup[lstlisting]{
  font=small,
  labelfont=bf,
  labelsep=period,
}

% ══════════════════════════════════════════════════════════
% BLOCKQUOTES: tcolorbox
% ══════════════════════════════════════════════════════════
\newtcolorbox{blockquote}{
  enhanced,
  colback=quoteBg,
  colframe=quoteBorder,
  boxrule=0pt,
  leftrule=3pt,
  arc=0pt,
  outer arc=0pt,
  left=14pt,
  right=10pt,
  top=10pt,
  bottom=10pt,
  fontupper=\itshape\color{quoteText}\small,
  before upper={\tikz[baseline=-0.5ex]{
    \node[font=\Large\color{quoteBorder}] at (0,0) {``};
  }\hspace{4pt}},
  after upper={\hspace{4pt}\tikz[baseline=-0.5ex]{
    \node[font=\Large\color{quoteBorder}] at (0,0) {''};
  }},
}

% ══════════════════════════════════════════════════════════
% HEADERS / FOOTERS
% ══════════════════════════════════════════════════════════
\pagestyle{fancy}
\fancyhf{}
\fancyhead[L]{%
  {\fontsize{8}{10}\selectfont\color{titleMid}\leftmark}}
\fancyhead[R]{%
  {\fontsize{8}{10}\selectfont\color{titleMid}%
    Document Element Templates}}
\fancyfoot[C]{%
  {\fontsize{8}{10}\selectfont\color{titleMid}%
    \rule{2cm}{0.3pt}\hspace{6pt}\thepage\hspace{6pt}\rule{2cm}{0.3pt}}}
\renewcommand{\headrulewidth}{0pt}
\renewcommand{\footrulewidth}{0pt}

% ─── Begin document ───
\begin{document}

\thispagestyle{empty}
\begin{tikzpicture}[remember picture, overlay]
  \fill[titleDark]
    (current page.north west)
    rectangle ([yshift=-6pt]current page.north east);
  \fill[titleRule]
    ([xshift=50pt]current page.north west)
    rectangle ([xshift=54pt]current page.south west);
\end{tikzpicture}

\vspace{18pt}
{\fontsize{24}{28}\selectfont\bfseries\color{titleDark}
  \hspace{22pt}Document Elements}\\[12pt]
{\fontsize{11}{13}\selectfont\color{titleMid}
  \hspace{22pt}Tables, Lists, Code, Blockquotes \& Figures}

\vspace{20pt}

% ══════════════════════════════════════════════════════════
% SECTION 1: TABLES
% ══════════════════════════════════════════════════════════
\section{Tables}

\subsection{Booktabs Table with Alternating Rows}
A clean, booktabs-styled table with dark header and alternating row colors
for improved readability in longer data sets.

\vspace{4pt}
\begin{table}[h]
\centering
\caption{Comparison of Consensus Algorithms}
\begin{tabular}{@{} l c c c @{}}
\toprule
\rowcolorhead
\textcolor{tableHeadText}{\textbf{Algorithm}} &
\textcolor{tableHeadText}{\textbf{Consensus}} &
\textcolor{tableHeadText}{\textbf{Tolerance}} &
\textcolor{tableHeadText}{\textbf{Complexity}} \\
\midrule
Paxos        & Safety    & $f < n/2$       & Moderate \\
\rowcoloralt
Raft         & Safety    & $f < n/2$       & Low      \\
Viewstamped  & Safety    & $f < n/2$       & Moderate \\
\rowcoloralt
PBFT         & Safety    & $f < n/3$       & High     \\
Zab          & Total     & $f < n/2$       & Moderate \\
\rowcoloralt
Tendermint   & Total     & $f < n/3$       & Moderate \\
HotStuff     & Total     & $f < n/3$       & Low      \\
\bottomrule
\end{tabular}
\end{table}

\subsection{Multi-Column Table with Merged Cells}
Tables can span multiple columns and merge cells for grouped data presentation
using the \texttt{multirow} package alongside booktabs.

\vspace{4pt}
\begin{table}[h]
\centering
\caption{Network Protocols by Layer}
\begin{tabular}{@{} l l l c @{}}
\toprule
\rowcolorhead
\textcolor{tableHeadText}{\textbf{Layer}} &
\textcolor{tableHeadText}{\textbf{Protocol}} &
\textcolor{tableHeadText}{\textbf{Transport}} &
\textcolor{tableHeadText}{\textbf{Reliable}} \\
\midrule
\multirow{3}{*}{Application} & HTTP  & TCP & Yes \\
\rowcoloralt
                           & gRPC  & HTTP/2 & Yes \\
                           & DNS   & UDP & No  \\
\midrule
\multirow{2}{*}{Transport}  & TCP   & ---  & Yes \\
\rowcoloralt
                           & QUIC  & UDP  & Yes \\
\midrule
\rowcoloralt
\multirow{2}{*}{Network}    & IPv4  & ---  & No  \\
                           & IPv6  & ---  & No  \\
\bottomrule
\end{tabular}
\end{table}

% ══════════════════════════════════════════════════════════
% SECTION 2: LISTS
% ══════════════════════════════════════════════════════════
\section{Lists}

\subsection{Itemize (Unordered)}
Custom-styled bullet points with three nesting levels, each using
a distinct TikZ-drawn marker for visual hierarchy.

\begin{itemize}
  \item Fault tolerance ensures system availability during failures
    \begin{itemize}
      \item Redundancy is the primary mechanism
        \begin{itemize}
          \item Active-passive and active-active configurations
          \item Geographic distribution across data centers
        \end{itemize}
      \item Health monitoring provides early failure detection
      \item Automatic failover minimizes downtime
    \end{itemize}
  \item Consistency models define data visibility guarantees
  \item Scalability enables handling growing workloads
    \begin{itemize}
      \item Horizontal scaling adds more machines
      \item Vertical scaling increases machine capacity
    \end{itemize}
\end{itemize}

\subsection{Enumerate (Ordered)}
Numbered lists with styled badges, useful for sequential processes
or ranked information presentation.

\begin{enumerate}
  \item Client sends request to the leader node
    \begin{enumerate}
      \item Leader appends entry to its local log
      \item Leader increments log index and term
    \end{enumerate}
  \item Leader replicates entry to follower nodes
  \item Followers acknowledge successful append
  \item Leader commits entry once majority acknowledges
  \item Leader notifies client of successful commit
\end{enumerate}

\subsection{Description Lists}
Term-definition pairs with bold, colored terms for quick scanning.
Ideal for glossaries, feature lists, or configuration options.

\begin{description}
  \item[Consensus] Agreement among distributed nodes on a single
    data value or system state, despite failures and network delays.
  \item[Quorum] The minimum number of nodes required to agree before
    an operation is considered successful, typically a majority.
  \item[Leader Election] The process by which a distributed system
    selects a coordinator node to manage replication and coordination.
  \item[Heartbeat] A periodic message exchanged between nodes to
    signal liveness and detect failures within bounded time.
  \item[Partition Tolerance] The ability of a system to continue
    operating despite network partitions that isolate subsets of nodes.
\end{description}

% ══════════════════════════════════════════════════════════
% SECTION 3: CODE LISTINGS
% ══════════════════════════════════════════════════════════
\section{Code Listings}

\subsection{Syntax Highlighting}
Code blocks with language-aware syntax highlighting, line numbers,
and a subtle left accent bar. Built on the \texttt{listings} package.

\begin{lstlisting}[language=Python, caption={Raft-style consensus implementation}]
class RaftNode:
    def __init__(self, node_id: int, peers: list):
        self.id = node_id
        self.peers = peers
        self.state = "follower"
        self.term = 0
        self.log = []
        self.commit_index = 0
        self.voted_for = None  # None means no vote cast

    def request_vote(self, candidate_term: int):
        """Handle incoming RequestVote RPC."""
        if candidate_term > self.term:
            self.term = candidate_term
            self.voted_for = None
            self.state = "follower"
        return self._grant_vote(candidate_term)

    def append_entries(self, leader_term: int):
        """Handle incoming AppendEntries (heartbeat)."""
        if leader_term >= self.term:
            self.term = leader_term
            self.state = "follower"
            return True  # Accept entry
        return False  # Reject stale leader
\end{lstlisting}

\subsection{Configuration File Example}
Inline code with monospaced font is also styled consistently.
Configuration values use \texttt{key = value} syntax.
The \texttt{\textbackslash texttt\{\}} command provides inline
monospaced text.

\begin{lstlisting}[language={}, caption={Cluster configuration file}]
# raft-cluster.toml
[cluster]
node_id = "node-1"
data_dir = "/var/lib/raft"
bind_addr = "0.0.0.0:7946"

[replication]
snapshot_interval = "5m"
max_log_entries = 65536
heartbeat_timeout = "500ms"
election_timeout = "1.5s"

[storage]
backend = "boltdb"
wal_dir = "/var/lib/raft/wal"
snapshot_dir = "/var/lib/raft/snap"
\end{lstlisting}

% ══════════════════════════════════════════════════════════
% SECTION 4: BLOCKQUOTES
% ══════════════════════════════════════════════════════════
\section{Blockquotes}

Styled blockquotes with a left accent bar and warm background,
suitable for citations, callouts, or important remarks.

\begin{blockquote}
There are only two hard things in Computer Science: cache invalidation
and naming things. --- Phil Karlton
\end{blockquote}

\vspace{6pt}

\begin{blockquote}
A distributed system is one in which the failure of a computer you
didn't even know existed can render your own computer unusable.
--- Leslie Lamport
\end{blockquote}

\vspace{6pt}

\begin{blockquote}
Any sufficiently complicated C or Fortran program contains an ad hoc,
informally-specified, bug-ridden, slow implementation of half of
Common Lisp. --- Greenspun's Tenth Rule
\end{blockquote}

% ══════════════════════════════════════════════════════════
% SECTION 5: FIGURES (placeholder)
% ══════════════════════════════════════════════════════════
\section{Figures}

\subsection{Figure with Caption}
A figure environment with a styled caption. Since no image files are
available in this demo, we use a TikZ-drawn placeholder.

\begin{figure}[h]
\centering
\begin{tikzpicture}
  \fill[codeBg] (0,0) rectangle (8,3.5);
  \draw[codeBorder, line width=0.6pt, rounded corners=3pt]
    (0,0) rectangle (8,3.5);
  \node[font=\footnotesize\ttfamily\color{titleMid}] at (4,3.1)
    {figure-placeholder.png};
  % Simulated architecture diagram
  \node[draw=titleRule, fill=white, rounded corners=2pt,
        minimum width=1.4cm, minimum height=0.7cm,
        font=\footnotesize\bfseries\color{titleDark}] (c) at (2,1.75) {Client};
  \node[draw=titleRule, fill=white, rounded corners=2pt,
        minimum width=1.4cm, minimum height=0.7cm,
        font=\footnotesize\bfseries\color{titleDark}] (l) at (4,1.75) {Leader};
  \node[draw=titleRule, fill=white, rounded corners=2pt,
        minimum width=1.4cm, minimum height=0.7cm,
        font=\footnotesize\bfseries\color{titleDark}] (f1) at (6,2.5) {F1};
  \node[draw=titleRule, fill=white, rounded corners=2pt,
        minimum width=1.4cm, minimum height=0.7cm,
        font=\footnotesize\bfseries\color{titleDark}] (f2) at (6,1.75) {F2};
  \node[draw=titleRule, fill=white, rounded corners=2pt,
        minimum width=1.4cm, minimum height=0.7cm,
        font=\footnotesize\bfseries\color{titleDark}] (f3) at (6,1.0) {F3};
  \draw[->, titleRule, thick] (c) -- (l);
  \draw[->, titleRule, thick] (l) -- (f1);
  \draw[->, titleRule, thick] (l) -- (f2);
  \draw[->, titleRule, thick] (l) -- (f3);
  \node[font=\tiny\color{titleMid}] at (4,0.4) {Raft Cluster Topology};
\end{tikzpicture}
\caption{A typical Raft cluster with one leader and three followers
  handling client requests via log replication.}
\end{figure}

\subsection{Mathematical Content}
Mathematical expressions and equations are rendered with consistent
styling using the \texttt{amsmath} package.

The CAP theorem states that a distributed data store can simultaneously
provide at most two of the following three guarantees:

\begin{description}
  \item[Consistency] Every read receives the most recent write or an error
  \item[Availability] Every request receives a response (not an error)
  \item[Partition Tolerance] System continues despite network partitions
\end{description}

The probability of quorum agreement in a system of $n$ nodes
tolerating $f$ failures is expressed as:

\vspace{4pt}
\[
  Q(n, f) = \binom{n}{f+1} \cdot p^{f+1} \cdot (1-p)^{n-f-1}
\]

where $p$ is the probability that any single node responds within
the timeout window and $f$ is the maximum number of tolerated failures
in a majority-based quorum system.

\end{document}
